\section{Methods}
\label{sec:methods}
We perform preliminary experiments to gauge the proficiency and behavior of current large language models on interactive coding tasks with Bash and SQL. To observe and elicit relevant reasoning skills, we draw on several existing prompting strategies that have been put forth to augment language models' reasoning and problem-solving skills. We apply these prompting strategies to models across the following three families: OpenAI (\texttt{text-davinci-003}, \texttt{gpt-3.5-turbo}, \texttt{gpt-4}), PaLM-2 (\texttt{text-bison-001}, \texttt{chat-bison-001}) \cite{anil2023palm}, and Open Source (Vicuna-13B \cite{vicuna2023}, StarChat-16B \cite{li2023starcoder}).

Figure~\ref{fig:prompting-approaches} visualizes the four adjusted prompting strategies we evaluate on \benchmark{}.

\textbf{Single Turn} is a zero-shot attempt. A model is given a simple description of the task setting and asked to generate code in a specific programming language that would address the query. The first generation in response to the user's question is then evaluated in the \benchmark{} environment.

\begin{figure}
    \centering
    \includegraphics[width=1.0\textwidth]{figures/assets/Approaches.pdf}
    \caption{Overview of Prompting Strategies adjusted for evaluation on \benchmark{}. The "Try Again" termination constraint is conditioned on reward = 1, while ReAct~\cite{yao2023react} and Plan \& Solve~\cite{wang2023planandsolve} are determined by the agent itself. This is because the purpose of the "Try Again" method is to explore how capable agents are at error correction from feedback, while the other two are more concerned with the overall success of general problem-solving strategies.}
    \label{fig:prompting-approaches}
\end{figure}

\textbf{"Try Again"} is an iterative feedback set up. In the initial message, the agent is informed of the task setting and its interactive nature; an agent has multiple turns to interact with the system, wherein each turn, upon generating an action, the execution output of the action is fed back as an observation. This continues until a reward of 1 (task completion) is achieved or the number of turns ($n$) is exhausted. The agent's position in this approach is meant to mirror human software development as closely as possible. The goal of this method is to probe language models' raw interactive coding abilities in addition to illustrating the benefits and different challenges that arise in interactive coding tasks.

\textbf{ReAct and Plan \& Solve.} We write prompts and design workflows that follow the text and task configurations described in ReAct~\cite{yao2023react} and Plan \& Solve~\cite{wang2023planandsolve} as faithfully as possible. For these two approaches, the termination of a task episode is conditioned upon the agent's own judgment, as our goal with these methods is to gauge the transferability to and efficacy of existing reasoning frameworks with respect to the interactive coding task.
Full prompt templates are included in \S\ref{appx:experiments:templates}.
